\documentclass[10pt,a4paper]{article}

\usepackage[margin=1.5cm]{geometry}
\usepackage{graphicx}
\usepackage[hidelinks]{hyperref}
\usepackage{caption}

\title{Music Theory Notes}
\author{}
\date{\today}

\setlength{\parskip}{0.3em}
\setlength{\textfloatsep}{5pt}
\setlength{\intextsep}{4pt}
\captionsetup{font=small,skip=2pt}

\begin{document}
\maketitle

\section{Example score: C major scale}

The score is generated from \texttt{lilypond\_files/example.ly}.
\begin{center}
  \includegraphics[width=0.72\textwidth,height=2.5cm,keepaspectratio]{score_pdfs/example.pdf}
  \captionof{figure}{C major scale.}
\end{center}
\noindent\textbf{Audio:} \href{run:audio/example.wav}{Listen to the C major scale (WAV)}.

\section{Modulation example: C major to D major}

This compact piano excerpt illustrates a modulation from C major to D major.
The harmony is \(\mathrm{C \;-\; C^+/B \;-\; Am \;-\; A^+ \;-\; D}\).
The score uses a piano grand staff and shows the change from C major to D major,
including the two-sharp key signature.

\begin{center}
  \includegraphics[width=0.92\textwidth,height=4.2cm,keepaspectratio]{score_pdfs/modulation.pdf}
  \captionof{figure}{Piano grand staff: modulation from C major to D major.}
\end{center}
\noindent\textbf{Audio:} \href{run:audio/modulation.wav}{Listen to the modulation example (WAV)}.

\end{document}
